 \chapterpage{1}{Estado del Arte}
\clearpage
\addtocontents{toc}{\protect\cftpagenumberson{chapter}}
\setcounter{page}{10} % O el número de página correcto
\refstepcounter{chapter} % incrementa el contador del capítulo
\addcontentsline{toc}{chapter}{Cap. 1 Estado del arte} % opcional: que sí salga en el índice
En este capitulo se va a ver como las tecnologías han estado evolucionando en los sistemas de gestión de ventas e inventarios en las pequeñas y medianas empresas, el que se ha hecho en cada época. También un análisis profundo de proyectos similares en universidades de México de 5 años para atrás, esto para darnos una idea mejor en lo que conlleva nuestro proyecto.
\section {Evolución de las tecnologías para los sistemas de gestión de ventas e inventarios en las pequeñas y medianas empresas.}



Es una historia de democratización tecnológica. Lo que antes era exclusivo de grandes corporaciones multinacionales (como el control en tiempo real) hoy en día esta al alcance de una tienda local gracias a la nube.

Esta transformación se puede dividir en cuatro grandes eras tecnológicas:
    \subsection{ Era de digitalización temprana: Hoja de cálculo y códigos de barra (años 90 - 2000).}

    Con la llegada de las computadoras personales, las PyMEs dieron su primer salto.

    La revolución de Excel, las hojas de calculo comenzaron a reemplazar al papel, ya que permiten cálculos automáticos, pero seguían siendo estáticas, si se vendía algo, se tenia que ir a la computadora a restarlo manualmente.
    
    El código de barras fue la primera gran tecnología de hardware que adoptaron las PyMEs, permitió agilizar el cobro y reducir errores humanos al teclear los precios.
    
    \subsection{Revolución de la nube (2010 - 2020).}

    Este es el punto de inflexión mas critico para las PyMEs. La tecnología cambio de modelo, ya no comprabas el software (como un CD en una caja), si no que lo rentabas por una suscripción mensual en las plataformas de streaming o música.

    El acceso remoto por primera vez, el dueño de la empresa podía ver sus ventas en tiempo real desde su celular estando en su casa. El sistema potente de planificación de recursos empresariales que costaban miles de dólares se volvieron accesibles por cuotas mensuales bajas.
    
    \subsection{Era de la inteligencia (2020 - Actualidad).}

    Hoy en día, las tecnologías de gestión ya no solo registran lo que paso, sino que predicen lo que pasara.

    La inteligencia artificial, sistemas modernos analizan patrones históricos haciendo que te manden mensajes de alerta para que no se te olvide rellenar stock en las diferentes temporadas del año. El IoT que usa etiquetas RFID que permiten contar todo el inventario de una bodega en segundos sin linea de vista directa.

\section{Proyectos Similares.}
\subsection{"Sistema de Información para el Control de Inventarios del Almacén del ITS"}
11 de Abril 2011, Instituto Tecnológico de Saltillo, Saltillo, México.

El proyecto está dirigido a hacer mejoras en el Instituto Tecnológico de Saltillo (ITS), en el Departamento de Recursos Materiales y Servicios, en el área de Almacén, como un soporte de apoyo administrativo que permita llevar el control de inventarios del almacén del ITS con la finalidad de ofrecer rapidez y seguridad en el manejo del inventario. Actualmente se trabaja de forma manual con formatos en Word y Excel, es necesario emigrar a otros entornos que permitan automatizar las tareas propias del almacén. Dentro de los resultados que se propondrán del proyecto, es obtener consultas e informes requeridos por los usuarios, un seguimiento y control de los materiales que entran y salen del almacén a los diferentes departamentos del ITS. 

El uso de un software para el control de inventarios es una herramienta que facilitará el proceso en las actividades de los usuarios del almacén. Dentro de este proyecto los beneficiados con esta investigación serán principalmente el personal que ahí labora (usuarios), quienes otorgan atención y servicio de abastecimiento de mercancía tanto de artículos de oficina, como equipo de cómputo, material especializado de laboratorios, de limpieza, entre otros a los diferentes departamentos del Instituto, los cuales contarán con un servicio de atención más rápida, eficiente, oportuna para cubrir sus necesidades básicas. Con el presente proyecto se pretende cubrir las necesidades del almacén del ITS con una herramienta enfocada a un Sistema de información para tener organizada la información concerniente al manejo de actividades básicas de Inventarios acordes a las necesidades de la institución, se considera que la mayoría de los Tecnológicos del país ya cuentan con sistema propio, nuestra institución no tiene un sistema de esa naturaleza y en base a este desarrollo se pretende cubrir esa necesidad a través de un sistema de control de inventarios.
Este control se lleva mediante tarjetas denominadas Kardex, en donde se lleva el registro de cada unidad, su valor de compra, la fecha de adquisición, el valor de la salida de cada unidad y la fecha en que se retira del inventario. De esta forma, en todo momento se puede conocer el saldo exacto de los inventarios y el valor del costo de venta, el control permanente de los sistemas en base a los inventarios existentes.

Inicialmente se realizó un análisis de requerimientos y se usaron herramientas de modelado tales como el diagrama de flujo de datos para representar los datos relevantes del diseño del sistema, el modelo incremental para plantear las etapas del desarrollo del proyecto. Se usó el modelo de desarrollo de prototipos para el desarrollo del software, el cual inicia con la recolección de requisitos y en donde el desarrollador y el usuario definen los objetivos globales para el software y el diseño se centra en una representación rápida de las tareas del sistema, el cual se va refinando en base a las necesidades de los usuarios.
Para el desarrollo de este proyecto se utilizó la metodología de programación extrema (PE). Esta metodología consta de un conjunto de reglas y prácticas comprendidas en cuatro actividades básicas; Planeación, diseño, codificación y pruebas. Inicialmente se conformó la historia del usuario, la planificación del proyecto y la localización del mismo, se utilizó la técnica de construcción del prototipo para realizar las interfases del sistema. Se usó la metodología incremental para hacer las diferentes incrementos del sistema, previamente aprobados por el usuario(s), se codificaron los módulos en visual Basic 6.0, se diseño la base de datos en Access 2003, contemplando sus tablas relacionales con sus respectivas relaciones y se enlazaron a Visual Basic para construir las interfaces del sistema, así como las pruebas pertinentes, desde la etapa de desarrollo hasta la etapa final de la aplicación.


\subsection{"SISTEMA CONTROL DE INVENTARIO"}
2018, Instituto Universitario Tecnológico de Antioquia, Medellín, Colombia.

La realización de este trabajo de investigación, se hizo con el propósito de mostrar las falencias que tiene el almacén de la fundación cementerio de san pedro, con la idea de mejorar dando respuesta y soluciones a muchos problemas por los cuales está pasando el inventario, no generando gastos en productos que se encuentran registrados en el almacén, debido a no contar con un esquema apropiado que vuelvan a generar las mismas órdenes de compra, empleando un método correcto con un sistema de stock de seguridad, las ordenes de compras corresponden a un modelo de gestión de inventario que controle los procesos, la intervención en la gestión de compras a proveedores, de materiales dando solución a los problemas actuales controlando gastos y alcanzando una mejor organización manteniendo el nivel óptimo, durante años no se implemento medidas de gestión, políticas y normas para asegurar el bienestar tanto para las personas directas o indirectamente relacionadas con el almacén, que busca aliviar una carga impositiva para el mejoramiento del espacio como la reducción de costos y gastos reflejándose en la parte contable de la empresa, diseñar una propuesta que gestione el control de inventarios con la eficiencia necesaria para solucionar uno de los muchos inconvenientes por las cuales está pasando la empresa, la idea de esta investigación es mostrar el camino de otras empresas para que se apoyen y busquen solucionar problemas de inventarios en sus organizaciones.
Inicialmente se realizó un análisis de requerimientos y se usaron herramientas de modelado tales como el diagrama de flujo de datos para representar los datos relevantes del diseño del sistema, el modelo incremental para plantear las etapas del desarrollo del proyecto. Se usó el modelo de desarrollo de prototipos para el desarrollo del software, el cual inicia con la recolección de requisitos y en donde el desarrollador y el usuario definen los objetivos globales para el software y el diseño se centra en una representación rápida de las tareas del sistema, el cual se va refinando en base a las necesidades de los usuarios.
Para el desarrollo de este proyecto se utilizó la metodología de programación extrema (PE). Esta metodología consta de un conjunto de reglas y prácticas comprendidas en cuatro actividades básicas; Planeación, diseño, codificación y pruebas. Inicialmente se conformó la historia del usuario, la planificación del proyecto y la localización del mismo, se utilizó la técnica de construcción del prototipo para realizar las interfases del sistema. Se usó la metodología incremental para hacer las diferentes incrementos del sistema, previamente aprobados por el usuario(s), se codificaron los módulos en visual Basic 6.0, se diseño la base de datos en Access 2003, contemplando sus tablas relacionales con sus respectivas relaciones y se enlazaron a Visual Basic para construir las interfaces del sistema, así como las pruebas pertinentes, desde la etapa de desarrollo hasta la etapa final de la aplicación.

\subsection{"DISEÑO DE UN SISTEMA DE CONTROL DE INVENTARIO Y ORGANIZACIÓN DE LAS BODEGAS DE PRODUCTO TERMINADO DE LA EMPRESA ECUAESPUMAS-LAMITEX S.A"}
Marzo 2018, Universidad Politécnica Salesiana, Cuenca, Ecuador.

El presente Proyecto Técnico consiste en el Diseño de un Sistema de Control de Inventarios y Organización de las Bodegas de Producto Terminado de la Empresa ECUAESPUMAS-LAMITEX S.A., en la cual se detallan problemas con respecto a la administración de sus inventarios y en la organización de sus bodegas de producto terminado. El campo de estudio se enfoca a las siguientes áreas: corte de esponja y traslado (despacho), empaque, bodegas de producto terminado, despacho y administrativa Para su desarrollo se utilizó una investigación descriptiva, a fin de detallar el estado actual en el que opera la Empresa con respecto al campo de estudio, empleando técnicas y herramientas de investigación las cuales se detallan en el Marco Metodológico. 
La Propuesta de Mejora está basada en el Sistema de Inventarios Periódico y la Metodología de las 5s, que se basa en la clasificación, organización, limpieza, estandarización, y la autodisciplina, lo cual nos va a dar las pautas necesarias para mejorar las condiciones de trabajo, reducir los gastos de tiempo y energía, reducir los riesgos de accidentes o sanitarios, mejorar la calidad del trabajo, seiri (clasificación), seiton (organizar), seiso (limpieza), seiketsu (estandarizar) y shitsuke (disciplina).

El presente proyecto técnico requerirá de una investigacion del tipo descriptiva a fin de recolectar información precisa del estado actual de las Bodegas de Producto Terminado en la Empresa ECUAESPUMAS-LAMITEX S.A.. Se empleó la técnica de la observacion directa, para comprender como se realizan las actividades en las presentes áreas de estudio, está información fue recolectada mediante fotografías y apuntes.

Esta investigación no aplica el cálculo de una población y muestra ya que no es una investigación estadística, por ello solo se especifica las áreas de estudio las cuales son las siguientes: corte de esponja y traslado “despacho”, empaque, bodegas de producto terminado y despacho.

Técnicas de recolección de información
\begin{itemize}  
 \item Investigación bibliográfica:
Se obtendrá información bibliográfica, de documentos y archivos de la empresa y de bibliotecas particulares, lo cual servirá como guía para desarrollar correctamente el Proyecto de Titulación.
 
 \item Observación directa:
Esta técnica de investigación ayudará a comprender a profundidad el funcionamiento actual de las bodegas de producto terminado, su mantenimiento, administración y control de inventario, al realizar un recorrido por estas así como por las diferentes áreas de estudio. La observación directa permitirá la identificación de problemas existentes, evitando influir en las actividades del personal involucrado a fin de recolectar la información lo más cercana y directa en el proceso administrativo y operativo de las bodegas.
\end{itemize}

Técnicas de análisis de información
\begin{itemize}
     \item Análisis de inventarios ABC:
Este análisis permitirá conocer de manera acertada el movimiento del producto terminado en el periodo de un año a fin de clasificar los productos según su demanda y rotación, ayudando a disminuir recorridos largos para el despacho de los productos, manteniéndolos organizados según su tipo, medidas y otras características.
\item Árbol de problemas:
Esta técnica ayudará a organizar toda la información recolectada y se utilizará con la finalidad de ilustrar de forma precisa las causas y efectos de los problemas evidenciados en el control y manejo de las bodegas de producto terminado, a fin de proponer soluciones acertadas para mitigar dichos problemas.
\item Diagrama de Pareto:
Este diagrama será empleado para identificar los productos 80-20, tomando la información del análisis ABC para clasificar los inventarios en tres categorías:
•	Productos de alta rotación
•	Productos de mediana rotación
•	Productos de baja rotación, especiales y bajo pedido


\end{itemize}

\subsection{"Implementación de un sistema de control de inventarios por medio del método primeras entradas primeras salidas (PEPS) en la empresa Comercializadora de Porcinos Mirasol S.A. de C.V."}
5 de Junio de 2018, Universidad Abierta y a Distancia de México, México.

Es necesario el establecimiento de un sistema de inventarios que permita tener un control adecuado del producto y evitar las mermas por sobre inventario, lo cual traerá como impacto el máximo aprovechamiento de los recursos, se eliminará la posibilidad de robos hormiga y el desperdicio por mermas y por producto caducado, además se evitarán los gastos generados por el tratamiento que debe darse a la carne echada a perder.
El principal problema que aqueja actualmente a la empresa es la falta de control de inventarios en el área de almacén (cámaras de refrigeración), pues se presentan diferencias al momento de contar las existencias, se dificulta la elaboración de reportes y se presentan mermas y desperdicios provocados por producto caducado, lo cual trae como consecuencia la necesidad de sanitizar las cámaras, además de que la carne echada a perder debe depositarse en una fosa especial y es necesario agregar químicos para que se desintegre rápidamente y con ello minimizar el impacto ambiental. Este proyecto reducirá pérdidas en la empresa al minimizar los desperdicios de mercancía, facilitará el trabajo del personal operativo y administrativo y representará un desarrollo del conocimiento para el personal, además de que ayudará a fomentar una cultura en la empresa.

Para poder lograr el objetivo general fue necesario trabajar sobre objetivos específicos, los cuales comenzaron con la realización de un inventario inicial, el cual se hizo fácilmente con el apoyo del ahora encargado de almacén, el gran detalle se presentó al momento de comparar el inventario realizado contra las existencias que se tenían registradas, pues se presentó una diferencia significativa, cuestión que posterior a la aplicación del nuevo sistema de control de inventarios se corrigió por completo; de igual manera se calcularon los mínimos, máximos y punto de reorden del producto, los cuales resultaron sumamente útiles para prevenir el desabasto y para evitar la acumulación de mercancía.
El diseño de los procedimientos fue todo un reto, pues el hecho de explicar a detalle y de manera perfectamente entendible para todo el personal los procesos involucrados en el control de inventarios significó un gran esfuerzo que al final trajo un gran aprendizaje.
Capacitar al personal para el manejo del sistema de control de inventarios fue el momento más satisfactorio del proyecto, puesto que en este punto se comparte con las personas involucradas en el proceso el arduo trabajo realizado durante casi un año y que representa un documento de valor para la empresa.

\subsection{"Popuesta de un sistema de gestión de inventarios para el almacén de la empresa COVIM S.A de C.V"}
Noviembre del 2019, Universidad de Sotavento, Coatzacoalcos, México.

Desarrollar estrategias de controll y organizacion de inventarios en la empresa COVIM S.A de C.V para el manejo de los insumos, le permitira mejorar sus procesos e incrementar la calidad del servicio al cliente en la ejecución de los proyectos. 
Determinar el Sistema de Inventario, presentar a la empresa un plan de control y rotación de inventarios, que permita a sus directivos mejorar el proceso de ingreso y salida de la empresa, se inicia con el análisis del inventario para la empresa para luego dar diferentes propuestas de gestión de inventario.
Las propuestas son:
\begin{itemize}
\item Propuesta del Sistema de Gestión de Inventarios Basado en las 5’S Japonesas, SEIRI (CLASIFICACIÓN)
En esta primera fase, nos centraremos en identificar y separar los repuestos y materiales necesarios de los innecesarios. El objetivo de esta fase es contar con un área de trabajo en donde solo estén los artículos y herramientas necesarias. SEITON (ORGANIZACIÓN) la Universidad de Sotavento A.C., determina cuestiones externas e internas mediante Consiste en arreglar u ordenar los artículos de la bodega, de manera que sea fácil y rápido encontrarlos, utilizarlos y reponerlos. El objetivo de esta tarea es que exista un lugar para cada cosa, adecuado a las rutinas de trabajo, listos para utilizarse y con su debida señalización. SEISO (LIMPIEZA) una vez que el espacio de trabajo está despejado (seiri) y ordenado (seiton), es mucho más fácil limpiarlo (seiso). Consiste en identificar y eliminar las fuentes de suciedad, asegurando que todos los medios se encuentran siempre en perfecto estado operativo. El incumplimiento de la limpieza puede tener muchas consecuencias, provocando incluso anomalías o el mal funcionamiento de la maquinaria.
SEIKETSU (ESTANDARIZAR) consiste en distinguir fácilmente una situación normal de otra anormal, mediante normas sencillas y visibles para todos. El objetivo en esta tarea es desarrollar condiciones de trabajo que eviten el retroceso de las primeras 3 S.
Una vez implementada las tres primeras S, el responsable de la bodega tiene que estar pendiente que estos tres procesos se cumplan a cabalidad (clasificar, ordenar y limpiar), se tiene que crear estándares. SHITSUKE (AUTODISCIPLINA) crear conciencia al personal sobre el orden y la limpieza, empezando por dar el ejemplo desde la gerencia. (No se le puede pedir a un mecánico que sea ordenado si ve a su supervisor inmediato realizando labores fuera de los parámetros de la 5S).
\item Propuesta del sistema de gestión de inventarios basado en el análisis ABC

La propuesta consiste en obtener un manejo de inventario más estricto por medio de la priorización de materiales, con base al método de control de inventarios ABC.
Se propone utilizar el método de control de inventario ABC o método de clasificación ABC para dar prioridad en cuanto a cantidad a solicitar y mantener en inventario desde el punto de vista monetario, es decir de mayor a menor costo:
•	A: Alto volumen monetario.
 
•	B: Volumen monetario medio.

•	C: Bajo volumen monetario.
\item Propuesta para la administración y control del inventario

La empresa no cuenta con la información necesaria de inventarios por lo que es de suma importancia proponer que se implementen algunos documentos que son importantes para poder administrar y controlar el inventario.
El hecho de controlar el inventario de manera eficaz es importante porque puede satisfacer las demandas de sus clientes con mayor rapidez.
La administración de inventarios tiene como meta, conciliar o equilibrar los siguientes objetivos:
•	Maximizar el servicio al cliente.

•	Maximizar la eficiencia de las unidades de compra, y

•	Minimizar la inversión en inventarios.

\item Toma Física del Inventario

La toma del inventario consiste en contar físicamente cada uno de los materiales, con el fin de dar el dato preciso o exacto de las cantidades de insumos en existencia.
En el conteo de un inventario resulta complejo verificar la exactitud de los datos con respecto al físico tomado anteriormente, se podría utilizar varios métodos para determinar las cantidades físicas en existencia, muchos de los cuales requieren la utilización de un equipo de trabajo.
\item Creación y Control del Banco de Datos

Una vez realizada la toma física del inventario, se procederá a crear la base de datos, para la cual se hace necesaria la existencia de un sistema en el cual se pueda registrar el inventario que se llevó a cabo con anterioridad.
El propósito del control y documentación de los materiales es de utilizar correctamente los sistemas de información para resolver la problemática que se enfrenta en el manejo de inventarios. Se propone implementar un sistema para el manejo de los registros referentes a los ingresos, egresos y devoluciones de materiales, de tal manera que el inventario físico coincida con los registros del sistema de control de inventarios. La información que dé como resultado de esta verificación permitirá comprar los materiales, repuestos o maquinarias sólo cuando se los requiera y adquirir sólo lo necesario
\item 3.4.3	Documentación y Control del Inventario

La empresa debe establecer una estrategia de control y verificación periódica de sus stocks de inventario, para eso debe de implementar la documentación necesaria de todas las operaciones relacionadas con los mismos.
A continuación de propondrá algunos formatos de los documentos necesarios para el control del inventario.

•	Orden de Compra

•	Requisición de Materiales

•	Guías de Remisión

•	Guías de Remisión

•	Informe de devolución a los Proveedores.

•	Stock de inventario o Tarjeta de Kárdex.
\end{itemize}
Con la propuesta de administración de inventarios se podrá tener un control más amplio de las mercaderías. Los formatos propuestos tienen un diseño de fácil lectura y entendimiento.
Bajo este marco se hace claro que el objetivo de la gestión del inventario es lograr un equilibrio entre la calidad de servicio brindado a los clientes y la inversión económica necesaria para ello, esto se ve traducido en una inversión inmovilizada que supone unos recursos financieros.


\subsection{"PROPUESTA PARA LA ADMINISTRACIÓN DE INVENTARIOS BASADA EN TECNOLOGÍA DE CÓDIGO DE BARRAS"}
Mayo 2023, Universidad Autónoma del Estado de Hidalgo, Hidalgo, México.

El presente proyecto surge debido a la problemática que presentan algunas empresas del sector metal mecánico para mantener un control en el flujo de los movimientos de los materiales, así como en los costos que éstos generan. Se propone desarrollar una propuesta para la implementación de tecnología basada en código de barras, mediante un caso de estudio en una empresa del sector metal mecánico, buscando reducir los tiempos y recursos necesarios utilizados para la identificación y acomodo de materia prima.

Un código de barras es un código basado en la representación de un conjunto de líneas paralelas de distinto grosor y espaciado que en su conjunto contienen una determinada información, es decir, las barras y espacios del código representan pequeñas cadenas de caracteres (Espinal, 2010).
De este modo, el código de barras permite reconocer rápidamente un artículo de forma única, global y no ambigua en un punto de la logística y así lograr realizar inventario o consultar sus características asociadas.

Como propuesta de solución tecnológica para este problema y después de realizar un análisis con el área encargada de la empresa, se considera lo siguiente:
\begin{itemize}
    \item Escáner de uso industrial con identificación de códigos inateck de uso industrial con batería de litio de duración de 85 horas, modelo LT 435-76. Este permitirá identificar la información contenida en el código de barras, codificarla y emitir o proporcionar la información a las diferentes áreas

    \item Conectores de USB de 35 cm. Permitirán realizar la interconexión entre las terminales y los lectores para el suministro de la información.
    \item 4 cargadores con base inteligente: después de conectar con una computadora por el cable USB incluido en el paquete, la base puede recargar el escáner de código de barras y puede transmitir el resultado de escaneo en la computadora no se necesita cargador ni receptor extra, 8 baterías de polímero de litio. Permitirá tener en todo momento energía para los lectores con la finalidad que en cada turno no se generen pausas por falta de pila y eso tenga repercusión en el desarrollo de la actividad.
    \item Instalación de software de operación en 4 computadoras. Permitirá codificar la información obtenida por el scanner, con la finalidad de tener la información para ser procesada y realizar la impresión correspondiente de las actividades a realizar.
    \item Software SAI ERP compatible con modelo de escáner e identificación de modelo de proceso; así como manejo de módulo de MM (similar a producto SAP)
\end{itemize}
Al efectuar la implementación del sistema de código de barras y los cambios de los procesos operativos que se tuvieron que hacer en la gestión de inventarios, los cuales duraron 120 días, se pudo observar el impacto, concluyendo lo siguiente:
\begin{itemize}
    \item Con la implementación de una tecnología en la industria como es el sistema de código de barras, se alcanzó el objetivo general de la actual investigación, incorporando metodología y procesos renovados que van de la mano con la tecnología implementada, mejorando significativamente el acceso de material y la asignación de lugares.
    \item Se encuentra actualizada en línea y se puede realizar la trazabilidad de los equipos, productos y mercaderías sin ningún inconveniente, brindando un mejor servicio a los clientes internos y externos, mejorando los tiempos de respuesta en el abastecimiento de la materia prima a producción, esto reduce los tiempos de proceso mejorando la eficiencia del proceso productivo.
    \item Al instalar el sistema de código de barras se ha conseguido eliminar los procedimientos manuales evitando los errores de registros y también las demoras en ubicar los equipos, productos y mercadería dentro de los almacenes para la asignación del tipo de material y proyecto al que pertenece.
    \item El sistema ofrece información correctamente registrada y localizada dentro de los almacenes, lo que lleva un mejor control y seguimiento pudiendo identificar a los responsables operativos, por consiguiente, se han reducido los tiempos operativos por tener exactitud en los inventarios y conocer las existencias y ubicación de los materiales en las instalaciones de planta.
\end{itemize}
